%!TEX root = ../nauchka.tex

\section{Заключение}

\begin{frame}
\frametitle{Результаты работы}

\begin{enumerate}
    \item \textbf{Tensor Parallelism:} линейное снижение пиковой памяти ($\approx 2\times$ при $P{=}2$), но потеря точности bounds для глубоких моделей (IBP fallback)

    \item \textbf{FSDP:} побитово идентичные bounds, экономия baseline-памяти 80--90\%, пиковой 34--39\%

    \item \textbf{Автоматическое шардирование:} обе функции работают с произвольными \texttt{BoundedModule}, включая ONNX-модели VNN-COMP

    \item \textbf{Теоретический анализ:} формальная модель $M_{\text{peak}}^{\text{FSDP}} = h^2(d/P + 1 + C)$ подтверждена экспериментально
\end{enumerate}

\end{frame}

\begin{frame}
\frametitle{Направления дальнейшей работы}

\begin{block}{Полная верификация через FSDP}
FSDP не изменяет вычислительный граф $\Rightarrow$ совместим с $\beta$-CROWN и Branch-and-Bound. Требуется распараллеливание управляющей логики BaB.
\end{block}

\vspace{0.3cm}

\begin{block}{Расширения}
\begin{itemize}
    \item Поддержка свёрточных слоёв (\texttt{BoundConv})
    \item Масштабирование на $P = 4, 8$ GPU
    \item Оптимизация $\alpha$-CROWN с TP и FSDP
\end{itemize}
\end{block}

\end{frame}
